\section{Model}
\label{sec:model}

\subsection{Scopes, entitlements and tasks}
Let $D$ be a database instance over relations $\mathcal{R}$. A \emph{scope} $\sigma$ maps each relation
$R$ in its domain to a pair $(C^\sigma_R,\varphi^\sigma_R)$ of a column set $C^\sigma_R\subseteq
\mathrm{attr}(R)$ and a row predicate $\varphi^\sigma_R$ over $\mathrm{attr}(R)$ ($\top$ if absent). An
actor holding exactly $\sigma$ sees, for each $R$, the view $V_\sigma(R)=\pi_{C^\sigma_R}\,
\sigma_{\varphi^\sigma_R}(R)$. We write $\sigma\sqsubseteq\tau$ when $\mathrm{dom}\,\sigma\subseteq
\mathrm{dom}\,\tau$ and, for every $R$, $C^\sigma_R\subseteq C^\tau_R$ and
$\varphi^\sigma_R\Rightarrow\varphi^\tau_R$. The meet $\sigma\sqcap\tau$ intersects columns and conjoins
predicates.

The \emph{entitlement} $E(u)$ of user $u$ is the scope the database grants $u$ through column privileges and
row-level-security policies. A \emph{task} $t$ is a natural-language request. For evaluation it comes with
a gold query $q_t$. The \emph{task scope} $S(t)$ is the scope provisioned for an agent acting for $u$ on
$t$, and we require $S(t)\sqsubseteq E(u)$.

\subsection{Needed data, sufficiency and minimality}
\begin{definition}[Needed data]
$N(t)$ maps each relation $R$ referenced by $q_t$ to the columns of $R$ that $q_t$ references (after
star expansion) and to the disjunction, over occurrences of $R$ in $q_t$, of the conjunction of that
occurrence's single-relation, subquery-free \texttt{WHERE} and inner-join conjuncts.
\end{definition}
Conjuncts that cannot be attributed to a single relation are dropped, so $N(t)$ over-approximates the data
$q_t$ touches. This choice can only make the over-exposure of the baselines look \emph{smaller}.

\begin{definition}[Sufficiency]
$\sigma$ is sufficient for $t$ on $D$ if $q_t$ references only columns in $\sigma$ and
$q_t(V_\sigma(D))=q_t(D)$.
\end{definition}
We test sufficiency by executing $q_t$ under the compiled scope and comparing result sets. A sufficient
$\sigma$ is \emph{minimal} if no $\sigma'\sqsubset\sigma$ is sufficient. We never claim that an inferred
scope is minimal (Section~\ref{sec:complexity}).

\subsection{Over-exposure ratio}
For a reach scope $\rho$ and a task whose needed data is permitted ($N(t)\sqsubseteq E(u)$), let
$\mathrm{rows}_D(\sigma,R)=|\sigma_{\varphi^\sigma_R}(R)|$, evaluated under $E(u)$'s row predicate. For B0,
$\rho=E(u)$; for TSE, $\rho=S(t)\sqcap E(u)$. We define
\begin{align}
\mathrm{OER}_{\mathrm{col}} &= \frac{\sum_{R}|C^\rho_R|}{\sum_{R}|C^{N}_R|},\\
\mathrm{OER}_{\mathrm{row}} &= \frac{\sum_{R\in\mathrm{dom}N}\mathrm{rows}_D(\rho,R)}{\sum_{R\in\mathrm{dom}N}\max(1,\mathrm{rows}_D(N,R))},\\
\mathrm{OER}_{\mathrm{cell}} &= \frac{\sum_{R\in\mathrm{dom}\rho}|C^\rho_R|\,\mathrm{rows}_D(\rho,R)}{\sum_{R\in\mathrm{dom}N}|C^N_R|\,\max(1,\mathrm{rows}_D(N,R))}.
\end{align}
The column and cell numerators sum over every relation the reach $\rho$ contains, so reachable but unneeded tables count; the row ratio is taken over the relations of $N(t)$ only, because rows of unneeded tables are already counted at the cell level. An OER of one means the agent can reach exactly the needed data. For an insufficient scope the OER can fall
below one, so we always report OER together with the sufficiency rate.

\subsection{Widening}
A widening request $w=(C_w,\Delta\varphi_w,\textit{reason})$ asks for additional columns and/or the removal
of task row predicates. The broker grants only the in-entitlement part, $w\sqcap E(u)$; everything else is
rejected and logged. An approval policy decides on the in-bound part (Section~\ref{sec:widening}). After
approval $S'(t)=S(t)\sqcup(w\sqcap E(u))$, so $S'(t)\sqsubseteq E(u)$ still holds.

\subsection{Trust assumptions and the design guarantee}
\label{sec:assumptions}
\begin{itemize}
\item[\textbf{A1}] \emph{DBMS correctness.} PostgreSQL enforces column privileges, row-level-security
policies, view-owner privilege semantics and \texttt{security\_barrier} predicate ordering as documented~\cite{pgrls,pgsecbarrier,pggrant}.
\item[\textbf{A2}] \emph{Single credential.} The agent's only database credential is the ephemeral login
created for the task.
\item[\textbf{A3}] \emph{Trusted broker.} The broker and its credential are outside the adversary's
control. The inference model may be wrong; it need not be trusted for the guarantee below.
\item[\textbf{A4}] \emph{No bypass privileges.} Neither the agent role nor the user role has
\texttt{SUPERUSER}, \texttt{BYPASSRLS}, table ownership, \texttt{CREATEROLE}, or membership with the
\texttt{SET} option in any other role, and \texttt{PUBLIC} has no privilege on base tables.
\end{itemize}

\begin{remark}[Design guarantee; holds by construction under A1--A4, not an empirical result]
\label{rem:guarantee}
Under A1--A4, every cell the agent can read through any query on its connection lies in
$\mathrm{reach}_D\big(S'(t)\sqcap E(u)\big)$, where $S'(t)$ is $S(t)$ joined with the approved
widenings.
\end{remark}
\noindent\emph{Proof sketch.} The login holds privileges only through membership (with \texttt{INHERIT},
without \texttt{SET}) in a per-task role that holds \texttt{USAGE} on the task schema and \texttt{SELECT}
on its views (A2, A4). Base tables are unreadable to it. Each view is owned by the user role, so by A1 it is
evaluated with that role's column grants and row policy and yields
$\pi_{C^S_R\cap C^E_R}\,\sigma_{\varphi^S_R\wedge\varphi^E_R}(R)$. The \texttt{security\_barrier}
option evaluates the view predicate before any user-supplied predicate (A1). Writes and object creation
are impossible because the scope role holds only \texttt{SELECT} on its views and no \texttt{CREATE} or
\texttt{TEMPORARY} privilege; the read-only default and statement timeout set on the login are defence in depth
only, since a session can change them. Widening
changes the views only through the bound check. \hfill$\square$

The Remark is conditional and holds by construction; we do not count it as an experimental finding.
Section~\ref{sec:res-security} reports our attempts to break it, and Section~\ref{sec:discussion} lists
what it does not cover: inference from in-scope data, catalog metadata (table and column names in
\texttt{pg\_catalog} are readable by every PostgreSQL role), timing and other covert channels, misuse of
in-scope data, a malicious user, and a compromised broker.

\subsection{A note on complexity}
\label{sec:complexity}
Deciding whether a column scope suffices for an arbitrary query on \emph{all} instances amounts to deciding
query containment or rewriting under views. Containment is undecidable for full relational algebra and
NP-complete already for conjunctive queries~\cite{chandra1977optimal}, as is rewriting conjunctive queries
using views~\cite{levy1995answering,halevy2001answering}. Moreover, $q_t$ is not available when the scope
is inferred; only the natural-language task is. Inference is therefore a prediction problem. We
(i)~never claim minimality, (ii)~define sufficiency per instance and test it by execution, and (iii)~bound
the cost of errors with widening.
